\begin{theorem}[The Hausdorff reflection of the boundary family]\hypertarget{tag-032F}{}\label{push:src:D.2.3}
\StatementRecord{push:src:D.2.3}{proved}{032F}{\SeriesRef{push:src:D.3.1}, \SeriesRef{eq:d-dual-exact}, \SeriesRef{push:src:D.2.2}}

Let \(\mathcal T=(\Delta_h\times\mathbb C^2)/\Lambda\) at the fixed
boundary parameter \(c_*\), after shrinking \(\Delta_h\). Two
central points represented by \((\eta,\xi)\) and
\((\eta',\xi')\) are nonseparated precisely when
\[
 \operatorname{Re}(\xi'-\xi)\in G_a
       :=\tfrac12\mathbb Z+a\mathbb Z\quad(a\in\mathbb Q),
\]
and every pair is nonseparated when \(a\notin\mathbb Q\).
There is no nonseparation between distinct base fibres, or between
distinct points of a noncentral fibre. The Hausdorff reflection
\(\mathcal T^{\mathrm H}\to\Delta_h\) has the original complex tori
over \(h\ne0\), and its central fibre is
\[
 \mathbb R/G_a\simeq S^1\quad(a\in\mathbb Q),\qquad
 \{*\}\quad(a\notin\mathbb Q).
\]
If \(a=m/n\) is in lowest terms, \(n>0\), then
\(G_a=\operatorname{lcm}(2,n)^{-1}\mathbb Z\).
\end{theorem}


\begin{proof}
\begin{enumerate}[label=\arabic*.,leftmargin=*]
\item We give the full limiting-group argument in \ref{push:src:D.3.1}; its formula also
applies to the fixed boundary family with \(c=c_*\), \(h\ne0\).
Put \(d_0=\operatorname{lcm}(2,n)\) in the rational case. In the dual
lattice, the fixed integral covector whose central limit is
\(d_0\operatorname{Re}\xi\) has
\((p,q,r)=(0,0,d_0)\). The exact dual formula
\eqref{eq:d-dual-exact} below shows that it extends continuously, as a
real-linear functional \(\ell_h\) on \(\mathbb C^2\), to that limit:
its last coefficient is \(O(h^2)\), its other moving coefficients
are \(O(h)\). Therefore
\[
 (h,z)\longmapsto\exp(2\pi i\ell_h(z))
\]
is continuous, invariant under every period, and restricts at the
centre to \(\exp(2\pi i d_0\operatorname{Re}\xi)\).
Points with different indicated central classes are separated by this
map to the circle.

\item For the converse take a ray \(h=\epsilon\omega\),
\(\epsilon\downarrow0\), whose real and imaginary parts have
irrational ratio. More precisely choose \(\omega\) so that
\(p\operatorname{Re}\omega+q\operatorname{Im}\omega\ne0\)
for every nonzero pair \((p,q)\) in the countable half-integral
lattice in \ref{push:src:D.3.1}. The asymptotic dual formula shows that all such
covectors escape to infinity. For rational \(a\), the surviving dual
group is exactly \(\mathbb Z d_0\operatorname{Re}\xi\); for
irrational \(a\), it is zero. Pontryagin--Chabauty duality
\cite[Theorem~3.1, p.~2017; proof, p.~2018]{Cornulier2011} now gives the period-group limits
\[
 \{z:d_0\operatorname{Re}\xi(z)\in\mathbb Z\},\qquad
 \mathbb C^2,
\]
respectively. Every element of the indicated limit is the limit of
period vectors. Such a sequence supplies intersecting images
of every two prescribed central neighbourhoods, proving
nonseparation. This includes the quadratic-scale differences of
\ref{push:src:D.2.2}, not only first-order differences.

\item Collapse exactly these central classes and leave every other fibre
unchanged. The resulting quotient is Hausdorff. The base coordinate
separates different fibres. Within a noncentral fibre, continuous
functions supported over a compact subdisc away from zero separate
distinct points, because there the original action is properly
discontinuous and gives a smooth torus bundle. In the rational central
fibre the displayed circle-valued function separates different
classes. These functions descend through this quotient, so
it is Hausdorff. Every continuous map from \(\mathcal T\) to a
Hausdorff space identifies nonseparated pairs and therefore factors
uniquely and continuously through this quotient. This proves the
universal property, including its topology. The elementary greatest
common divisor calculation gives the stated value of \(G_a\).
\end{enumerate}
\end{proof}